\documentclass[11pt]{article}

    \usepackage[breakable]{tcolorbox}
    \usepackage{parskip} % Stop auto-indenting (to mimic markdown behaviour)
    

    % Basic figure setup, for now with no caption control since it's done
    % automatically by Pandoc (which extracts ![](path) syntax from Markdown).
    \usepackage{graphicx}
    % Keep aspect ratio if custom image width or height is specified
    \setkeys{Gin}{keepaspectratio}
    % Maintain compatibility with old templates. Remove in nbconvert 6.0
    \let\Oldincludegraphics\includegraphics
    % Ensure that by default, figures have no caption (until we provide a
    % proper Figure object with a Caption API and a way to capture that
    % in the conversion process - todo).
    \usepackage{caption}
    \DeclareCaptionFormat{nocaption}{}
    \captionsetup{format=nocaption,aboveskip=0pt,belowskip=0pt}

    \usepackage{float}
    \floatplacement{figure}{H} % forces figures to be placed at the correct location
    \usepackage{xcolor} % Allow colors to be defined
    \usepackage{enumerate} % Needed for markdown enumerations to work
    \usepackage{geometry} % Used to adjust the document margins
    \usepackage{amsmath} % Equations
    \usepackage{amssymb} % Equations
    \usepackage{textcomp} % defines textquotesingle
    % Hack from http://tex.stackexchange.com/a/47451/13684:
    \AtBeginDocument{%
        \def\PYZsq{\textquotesingle}% Upright quotes in Pygmentized code
    }
    \usepackage{upquote} % Upright quotes for verbatim code
    \usepackage{eurosym} % defines \euro

    \usepackage{iftex}
    \ifPDFTeX
        \usepackage[T1]{fontenc}
        \IfFileExists{alphabeta.sty}{
              \usepackage{alphabeta}
          }{
              \usepackage[mathletters]{ucs}
              \usepackage[utf8x]{inputenc}
          }
    \else
        \usepackage{fontspec}
        \usepackage{unicode-math}
    \fi

    \usepackage{fancyvrb} % verbatim replacement that allows latex
    \usepackage{grffile} % extends the file name processing of package graphics
                         % to support a larger range
    \makeatletter % fix for old versions of grffile with XeLaTeX
    \@ifpackagelater{grffile}{2019/11/01}
    {
      % Do nothing on new versions
    }
    {
      \def\Gread@@xetex#1{%
        \IfFileExists{"\Gin@base".bb}%
        {\Gread@eps{\Gin@base.bb}}%
        {\Gread@@xetex@aux#1}%
      }
    }
    \makeatother
    \usepackage[Export]{adjustbox} % Used to constrain images to a maximum size
    \adjustboxset{max size={0.9\linewidth}{0.9\paperheight}}

    % The hyperref package gives us a pdf with properly built
    % internal navigation ('pdf bookmarks' for the table of contents,
    % internal cross-reference links, web links for URLs, etc.)
    \usepackage{hyperref}
    % The default LaTeX title has an obnoxious amount of whitespace. By default,
    % titling removes some of it. It also provides customization options.
    \usepackage{titling}
    \usepackage{longtable} % longtable support required by pandoc >1.10
    \usepackage{booktabs}  % table support for pandoc > 1.12.2
    \usepackage{array}     % table support for pandoc >= 2.11.3
    \usepackage{calc}      % table minipage width calculation for pandoc >= 2.11.1
    \usepackage[inline]{enumitem} % IRkernel/repr support (it uses the enumerate* environment)
    \usepackage[normalem]{ulem} % ulem is needed to support strikethroughs (\sout)
                                % normalem makes italics be italics, not underlines
    \usepackage{soul}      % strikethrough (\st) support for pandoc >= 3.0.0
    \usepackage{mathrsfs}
    

    
    % Colors for the hyperref package
    \definecolor{urlcolor}{rgb}{0,.145,.698}
    \definecolor{linkcolor}{rgb}{.71,0.21,0.01}
    \definecolor{citecolor}{rgb}{.12,.54,.11}

    % ANSI colors
    \definecolor{ansi-black}{HTML}{3E424D}
    \definecolor{ansi-black-intense}{HTML}{282C36}
    \definecolor{ansi-red}{HTML}{E75C58}
    \definecolor{ansi-red-intense}{HTML}{B22B31}
    \definecolor{ansi-green}{HTML}{00A250}
    \definecolor{ansi-green-intense}{HTML}{007427}
    \definecolor{ansi-yellow}{HTML}{DDB62B}
    \definecolor{ansi-yellow-intense}{HTML}{B27D12}
    \definecolor{ansi-blue}{HTML}{208FFB}
    \definecolor{ansi-blue-intense}{HTML}{0065CA}
    \definecolor{ansi-magenta}{HTML}{D160C4}
    \definecolor{ansi-magenta-intense}{HTML}{A03196}
    \definecolor{ansi-cyan}{HTML}{60C6C8}
    \definecolor{ansi-cyan-intense}{HTML}{258F8F}
    \definecolor{ansi-white}{HTML}{C5C1B4}
    \definecolor{ansi-white-intense}{HTML}{A1A6B2}
    \definecolor{ansi-default-inverse-fg}{HTML}{FFFFFF}
    \definecolor{ansi-default-inverse-bg}{HTML}{000000}

    % common color for the border for error outputs.
    \definecolor{outerrorbackground}{HTML}{FFDFDF}

    % commands and environments needed by pandoc snippets
    % extracted from the output of `pandoc -s`
    \providecommand{\tightlist}{%
      \setlength{\itemsep}{0pt}\setlength{\parskip}{0pt}}
    \DefineVerbatimEnvironment{Highlighting}{Verbatim}{commandchars=\\\{\}}
    % Add ',fontsize=\small' for more characters per line
    \newenvironment{Shaded}{}{}
    \newcommand{\KeywordTok}[1]{\textcolor[rgb]{0.00,0.44,0.13}{\textbf{{#1}}}}
    \newcommand{\DataTypeTok}[1]{\textcolor[rgb]{0.56,0.13,0.00}{{#1}}}
    \newcommand{\DecValTok}[1]{\textcolor[rgb]{0.25,0.63,0.44}{{#1}}}
    \newcommand{\BaseNTok}[1]{\textcolor[rgb]{0.25,0.63,0.44}{{#1}}}
    \newcommand{\FloatTok}[1]{\textcolor[rgb]{0.25,0.63,0.44}{{#1}}}
    \newcommand{\CharTok}[1]{\textcolor[rgb]{0.25,0.44,0.63}{{#1}}}
    \newcommand{\StringTok}[1]{\textcolor[rgb]{0.25,0.44,0.63}{{#1}}}
    \newcommand{\CommentTok}[1]{\textcolor[rgb]{0.38,0.63,0.69}{\textit{{#1}}}}
    \newcommand{\OtherTok}[1]{\textcolor[rgb]{0.00,0.44,0.13}{{#1}}}
    \newcommand{\AlertTok}[1]{\textcolor[rgb]{1.00,0.00,0.00}{\textbf{{#1}}}}
    \newcommand{\FunctionTok}[1]{\textcolor[rgb]{0.02,0.16,0.49}{{#1}}}
    \newcommand{\RegionMarkerTok}[1]{{#1}}
    \newcommand{\ErrorTok}[1]{\textcolor[rgb]{1.00,0.00,0.00}{\textbf{{#1}}}}
    \newcommand{\NormalTok}[1]{{#1}}

    % Additional commands for more recent versions of Pandoc
    \newcommand{\ConstantTok}[1]{\textcolor[rgb]{0.53,0.00,0.00}{{#1}}}
    \newcommand{\SpecialCharTok}[1]{\textcolor[rgb]{0.25,0.44,0.63}{{#1}}}
    \newcommand{\VerbatimStringTok}[1]{\textcolor[rgb]{0.25,0.44,0.63}{{#1}}}
    \newcommand{\SpecialStringTok}[1]{\textcolor[rgb]{0.73,0.40,0.53}{{#1}}}
    \newcommand{\ImportTok}[1]{{#1}}
    \newcommand{\DocumentationTok}[1]{\textcolor[rgb]{0.73,0.13,0.13}{\textit{{#1}}}}
    \newcommand{\AnnotationTok}[1]{\textcolor[rgb]{0.38,0.63,0.69}{\textbf{\textit{{#1}}}}}
    \newcommand{\CommentVarTok}[1]{\textcolor[rgb]{0.38,0.63,0.69}{\textbf{\textit{{#1}}}}}
    \newcommand{\VariableTok}[1]{\textcolor[rgb]{0.10,0.09,0.49}{{#1}}}
    \newcommand{\ControlFlowTok}[1]{\textcolor[rgb]{0.00,0.44,0.13}{\textbf{{#1}}}}
    \newcommand{\OperatorTok}[1]{\textcolor[rgb]{0.40,0.40,0.40}{{#1}}}
    \newcommand{\BuiltInTok}[1]{{#1}}
    \newcommand{\ExtensionTok}[1]{{#1}}
    \newcommand{\PreprocessorTok}[1]{\textcolor[rgb]{0.74,0.48,0.00}{{#1}}}
    \newcommand{\AttributeTok}[1]{\textcolor[rgb]{0.49,0.56,0.16}{{#1}}}
    \newcommand{\InformationTok}[1]{\textcolor[rgb]{0.38,0.63,0.69}{\textbf{\textit{{#1}}}}}
    \newcommand{\WarningTok}[1]{\textcolor[rgb]{0.38,0.63,0.69}{\textbf{\textit{{#1}}}}}


    % Define a nice break command that doesn't care if a line doesn't already
    % exist.
    \def\br{\hspace*{\fill} \\* }
    % Math Jax compatibility definitions
    \def\gt{>}
    \def\lt{<}
    \let\Oldtex\TeX
    \let\Oldlatex\LaTeX
    \renewcommand{\TeX}{\textrm{\Oldtex}}
    \renewcommand{\LaTeX}{\textrm{\Oldlatex}}
    % Document parameters
    % Document title
    \title{Unsupervised Machine Learning Project}
    
    
    
    
    
    
    
% Pygments definitions
\makeatletter
\def\PY@reset{\let\PY@it=\relax \let\PY@bf=\relax%
    \let\PY@ul=\relax \let\PY@tc=\relax%
    \let\PY@bc=\relax \let\PY@ff=\relax}
\def\PY@tok#1{\csname PY@tok@#1\endcsname}
\def\PY@toks#1+{\ifx\relax#1\empty\else%
    \PY@tok{#1}\expandafter\PY@toks\fi}
\def\PY@do#1{\PY@bc{\PY@tc{\PY@ul{%
    \PY@it{\PY@bf{\PY@ff{#1}}}}}}}
\def\PY#1#2{\PY@reset\PY@toks#1+\relax+\PY@do{#2}}

\@namedef{PY@tok@w}{\def\PY@tc##1{\textcolor[rgb]{0.73,0.73,0.73}{##1}}}
\@namedef{PY@tok@c}{\let\PY@it=\textit\def\PY@tc##1{\textcolor[rgb]{0.24,0.48,0.48}{##1}}}
\@namedef{PY@tok@cp}{\def\PY@tc##1{\textcolor[rgb]{0.61,0.40,0.00}{##1}}}
\@namedef{PY@tok@k}{\let\PY@bf=\textbf\def\PY@tc##1{\textcolor[rgb]{0.00,0.50,0.00}{##1}}}
\@namedef{PY@tok@kp}{\def\PY@tc##1{\textcolor[rgb]{0.00,0.50,0.00}{##1}}}
\@namedef{PY@tok@kt}{\def\PY@tc##1{\textcolor[rgb]{0.69,0.00,0.25}{##1}}}
\@namedef{PY@tok@o}{\def\PY@tc##1{\textcolor[rgb]{0.40,0.40,0.40}{##1}}}
\@namedef{PY@tok@ow}{\let\PY@bf=\textbf\def\PY@tc##1{\textcolor[rgb]{0.67,0.13,1.00}{##1}}}
\@namedef{PY@tok@nb}{\def\PY@tc##1{\textcolor[rgb]{0.00,0.50,0.00}{##1}}}
\@namedef{PY@tok@nf}{\def\PY@tc##1{\textcolor[rgb]{0.00,0.00,1.00}{##1}}}
\@namedef{PY@tok@nc}{\let\PY@bf=\textbf\def\PY@tc##1{\textcolor[rgb]{0.00,0.00,1.00}{##1}}}
\@namedef{PY@tok@nn}{\let\PY@bf=\textbf\def\PY@tc##1{\textcolor[rgb]{0.00,0.00,1.00}{##1}}}
\@namedef{PY@tok@ne}{\let\PY@bf=\textbf\def\PY@tc##1{\textcolor[rgb]{0.80,0.25,0.22}{##1}}}
\@namedef{PY@tok@nv}{\def\PY@tc##1{\textcolor[rgb]{0.10,0.09,0.49}{##1}}}
\@namedef{PY@tok@no}{\def\PY@tc##1{\textcolor[rgb]{0.53,0.00,0.00}{##1}}}
\@namedef{PY@tok@nl}{\def\PY@tc##1{\textcolor[rgb]{0.46,0.46,0.00}{##1}}}
\@namedef{PY@tok@ni}{\let\PY@bf=\textbf\def\PY@tc##1{\textcolor[rgb]{0.44,0.44,0.44}{##1}}}
\@namedef{PY@tok@na}{\def\PY@tc##1{\textcolor[rgb]{0.41,0.47,0.13}{##1}}}
\@namedef{PY@tok@nt}{\let\PY@bf=\textbf\def\PY@tc##1{\textcolor[rgb]{0.00,0.50,0.00}{##1}}}
\@namedef{PY@tok@nd}{\def\PY@tc##1{\textcolor[rgb]{0.67,0.13,1.00}{##1}}}
\@namedef{PY@tok@s}{\def\PY@tc##1{\textcolor[rgb]{0.73,0.13,0.13}{##1}}}
\@namedef{PY@tok@sd}{\let\PY@it=\textit\def\PY@tc##1{\textcolor[rgb]{0.73,0.13,0.13}{##1}}}
\@namedef{PY@tok@si}{\let\PY@bf=\textbf\def\PY@tc##1{\textcolor[rgb]{0.64,0.35,0.47}{##1}}}
\@namedef{PY@tok@se}{\let\PY@bf=\textbf\def\PY@tc##1{\textcolor[rgb]{0.67,0.36,0.12}{##1}}}
\@namedef{PY@tok@sr}{\def\PY@tc##1{\textcolor[rgb]{0.64,0.35,0.47}{##1}}}
\@namedef{PY@tok@ss}{\def\PY@tc##1{\textcolor[rgb]{0.10,0.09,0.49}{##1}}}
\@namedef{PY@tok@sx}{\def\PY@tc##1{\textcolor[rgb]{0.00,0.50,0.00}{##1}}}
\@namedef{PY@tok@m}{\def\PY@tc##1{\textcolor[rgb]{0.40,0.40,0.40}{##1}}}
\@namedef{PY@tok@gh}{\let\PY@bf=\textbf\def\PY@tc##1{\textcolor[rgb]{0.00,0.00,0.50}{##1}}}
\@namedef{PY@tok@gu}{\let\PY@bf=\textbf\def\PY@tc##1{\textcolor[rgb]{0.50,0.00,0.50}{##1}}}
\@namedef{PY@tok@gd}{\def\PY@tc##1{\textcolor[rgb]{0.63,0.00,0.00}{##1}}}
\@namedef{PY@tok@gi}{\def\PY@tc##1{\textcolor[rgb]{0.00,0.52,0.00}{##1}}}
\@namedef{PY@tok@gr}{\def\PY@tc##1{\textcolor[rgb]{0.89,0.00,0.00}{##1}}}
\@namedef{PY@tok@ge}{\let\PY@it=\textit}
\@namedef{PY@tok@gs}{\let\PY@bf=\textbf}
\@namedef{PY@tok@ges}{\let\PY@bf=\textbf\let\PY@it=\textit}
\@namedef{PY@tok@gp}{\let\PY@bf=\textbf\def\PY@tc##1{\textcolor[rgb]{0.00,0.00,0.50}{##1}}}
\@namedef{PY@tok@go}{\def\PY@tc##1{\textcolor[rgb]{0.44,0.44,0.44}{##1}}}
\@namedef{PY@tok@gt}{\def\PY@tc##1{\textcolor[rgb]{0.00,0.27,0.87}{##1}}}
\@namedef{PY@tok@err}{\def\PY@bc##1{{\setlength{\fboxsep}{\string -\fboxrule}\fcolorbox[rgb]{1.00,0.00,0.00}{1,1,1}{\strut ##1}}}}
\@namedef{PY@tok@kc}{\let\PY@bf=\textbf\def\PY@tc##1{\textcolor[rgb]{0.00,0.50,0.00}{##1}}}
\@namedef{PY@tok@kd}{\let\PY@bf=\textbf\def\PY@tc##1{\textcolor[rgb]{0.00,0.50,0.00}{##1}}}
\@namedef{PY@tok@kn}{\let\PY@bf=\textbf\def\PY@tc##1{\textcolor[rgb]{0.00,0.50,0.00}{##1}}}
\@namedef{PY@tok@kr}{\let\PY@bf=\textbf\def\PY@tc##1{\textcolor[rgb]{0.00,0.50,0.00}{##1}}}
\@namedef{PY@tok@bp}{\def\PY@tc##1{\textcolor[rgb]{0.00,0.50,0.00}{##1}}}
\@namedef{PY@tok@fm}{\def\PY@tc##1{\textcolor[rgb]{0.00,0.00,1.00}{##1}}}
\@namedef{PY@tok@vc}{\def\PY@tc##1{\textcolor[rgb]{0.10,0.09,0.49}{##1}}}
\@namedef{PY@tok@vg}{\def\PY@tc##1{\textcolor[rgb]{0.10,0.09,0.49}{##1}}}
\@namedef{PY@tok@vi}{\def\PY@tc##1{\textcolor[rgb]{0.10,0.09,0.49}{##1}}}
\@namedef{PY@tok@vm}{\def\PY@tc##1{\textcolor[rgb]{0.10,0.09,0.49}{##1}}}
\@namedef{PY@tok@sa}{\def\PY@tc##1{\textcolor[rgb]{0.73,0.13,0.13}{##1}}}
\@namedef{PY@tok@sb}{\def\PY@tc##1{\textcolor[rgb]{0.73,0.13,0.13}{##1}}}
\@namedef{PY@tok@sc}{\def\PY@tc##1{\textcolor[rgb]{0.73,0.13,0.13}{##1}}}
\@namedef{PY@tok@dl}{\def\PY@tc##1{\textcolor[rgb]{0.73,0.13,0.13}{##1}}}
\@namedef{PY@tok@s2}{\def\PY@tc##1{\textcolor[rgb]{0.73,0.13,0.13}{##1}}}
\@namedef{PY@tok@sh}{\def\PY@tc##1{\textcolor[rgb]{0.73,0.13,0.13}{##1}}}
\@namedef{PY@tok@s1}{\def\PY@tc##1{\textcolor[rgb]{0.73,0.13,0.13}{##1}}}
\@namedef{PY@tok@mb}{\def\PY@tc##1{\textcolor[rgb]{0.40,0.40,0.40}{##1}}}
\@namedef{PY@tok@mf}{\def\PY@tc##1{\textcolor[rgb]{0.40,0.40,0.40}{##1}}}
\@namedef{PY@tok@mh}{\def\PY@tc##1{\textcolor[rgb]{0.40,0.40,0.40}{##1}}}
\@namedef{PY@tok@mi}{\def\PY@tc##1{\textcolor[rgb]{0.40,0.40,0.40}{##1}}}
\@namedef{PY@tok@il}{\def\PY@tc##1{\textcolor[rgb]{0.40,0.40,0.40}{##1}}}
\@namedef{PY@tok@mo}{\def\PY@tc##1{\textcolor[rgb]{0.40,0.40,0.40}{##1}}}
\@namedef{PY@tok@ch}{\let\PY@it=\textit\def\PY@tc##1{\textcolor[rgb]{0.24,0.48,0.48}{##1}}}
\@namedef{PY@tok@cm}{\let\PY@it=\textit\def\PY@tc##1{\textcolor[rgb]{0.24,0.48,0.48}{##1}}}
\@namedef{PY@tok@cpf}{\let\PY@it=\textit\def\PY@tc##1{\textcolor[rgb]{0.24,0.48,0.48}{##1}}}
\@namedef{PY@tok@c1}{\let\PY@it=\textit\def\PY@tc##1{\textcolor[rgb]{0.24,0.48,0.48}{##1}}}
\@namedef{PY@tok@cs}{\let\PY@it=\textit\def\PY@tc##1{\textcolor[rgb]{0.24,0.48,0.48}{##1}}}

\def\PYZbs{\char`\\}
\def\PYZus{\char`\_}
\def\PYZob{\char`\{}
\def\PYZcb{\char`\}}
\def\PYZca{\char`\^}
\def\PYZam{\char`\&}
\def\PYZlt{\char`\<}
\def\PYZgt{\char`\>}
\def\PYZsh{\char`\#}
\def\PYZpc{\char`\%}
\def\PYZdl{\char`\$}
\def\PYZhy{\char`\-}
\def\PYZsq{\char`\'}
\def\PYZdq{\char`\"}
\def\PYZti{\char`\~}
% for compatibility with earlier versions
\def\PYZat{@}
\def\PYZlb{[}
\def\PYZrb{]}
\makeatother


    % For linebreaks inside Verbatim environment from package fancyvrb.
    \makeatletter
        \newbox\Wrappedcontinuationbox
        \newbox\Wrappedvisiblespacebox
        \newcommand*\Wrappedvisiblespace {\textcolor{red}{\textvisiblespace}}
        \newcommand*\Wrappedcontinuationsymbol {\textcolor{red}{\llap{\tiny$\m@th\hookrightarrow$}}}
        \newcommand*\Wrappedcontinuationindent {3ex }
        \newcommand*\Wrappedafterbreak {\kern\Wrappedcontinuationindent\copy\Wrappedcontinuationbox}
        % Take advantage of the already applied Pygments mark-up to insert
        % potential linebreaks for TeX processing.
        %        {, <, #, %, $, ' and ": go to next line.
        %        _, }, ^, &, >, - and ~: stay at end of broken line.
        % Use of \textquotesingle for straight quote.
        \newcommand*\Wrappedbreaksatspecials {%
            \def\PYGZus{\discretionary{\char`\_}{\Wrappedafterbreak}{\char`\_}}%
            \def\PYGZob{\discretionary{}{\Wrappedafterbreak\char`\{}{\char`\{}}%
            \def\PYGZcb{\discretionary{\char`\}}{\Wrappedafterbreak}{\char`\}}}%
            \def\PYGZca{\discretionary{\char`\^}{\Wrappedafterbreak}{\char`\^}}%
            \def\PYGZam{\discretionary{\char`\&}{\Wrappedafterbreak}{\char`\&}}%
            \def\PYGZlt{\discretionary{}{\Wrappedafterbreak\char`\<}{\char`\<}}%
            \def\PYGZgt{\discretionary{\char`\>}{\Wrappedafterbreak}{\char`\>}}%
            \def\PYGZsh{\discretionary{}{\Wrappedafterbreak\char`\#}{\char`\#}}%
            \def\PYGZpc{\discretionary{}{\Wrappedafterbreak\char`\%}{\char`\%}}%
            \def\PYGZdl{\discretionary{}{\Wrappedafterbreak\char`\$}{\char`\$}}%
            \def\PYGZhy{\discretionary{\char`\-}{\Wrappedafterbreak}{\char`\-}}%
            \def\PYGZsq{\discretionary{}{\Wrappedafterbreak\textquotesingle}{\textquotesingle}}%
            \def\PYGZdq{\discretionary{}{\Wrappedafterbreak\char`\"}{\char`\"}}%
            \def\PYGZti{\discretionary{\char`\~}{\Wrappedafterbreak}{\char`\~}}%
        }
        % Some characters . , ; ? ! / are not pygmentized.
        % This macro makes them "active" and they will insert potential linebreaks
        \newcommand*\Wrappedbreaksatpunct {%
            \lccode`\~`\.\lowercase{\def~}{\discretionary{\hbox{\char`\.}}{\Wrappedafterbreak}{\hbox{\char`\.}}}%
            \lccode`\~`\,\lowercase{\def~}{\discretionary{\hbox{\char`\,}}{\Wrappedafterbreak}{\hbox{\char`\,}}}%
            \lccode`\~`\;\lowercase{\def~}{\discretionary{\hbox{\char`\;}}{\Wrappedafterbreak}{\hbox{\char`\;}}}%
            \lccode`\~`\:\lowercase{\def~}{\discretionary{\hbox{\char`\:}}{\Wrappedafterbreak}{\hbox{\char`\:}}}%
            \lccode`\~`\?\lowercase{\def~}{\discretionary{\hbox{\char`\?}}{\Wrappedafterbreak}{\hbox{\char`\?}}}%
            \lccode`\~`\!\lowercase{\def~}{\discretionary{\hbox{\char`\!}}{\Wrappedafterbreak}{\hbox{\char`\!}}}%
            \lccode`\~`\/\lowercase{\def~}{\discretionary{\hbox{\char`\/}}{\Wrappedafterbreak}{\hbox{\char`\/}}}%
            \catcode`\.\active
            \catcode`\,\active
            \catcode`\;\active
            \catcode`\:\active
            \catcode`\?\active
            \catcode`\!\active
            \catcode`\/\active
            \lccode`\~`\~
        }
    \makeatother

    \let\OriginalVerbatim=\Verbatim
    \makeatletter
    \renewcommand{\Verbatim}[1][1]{%
        %\parskip\z@skip
        \sbox\Wrappedcontinuationbox {\Wrappedcontinuationsymbol}%
        \sbox\Wrappedvisiblespacebox {\FV@SetupFont\Wrappedvisiblespace}%
        \def\FancyVerbFormatLine ##1{\hsize\linewidth
            \vtop{\raggedright\hyphenpenalty\z@\exhyphenpenalty\z@
                \doublehyphendemerits\z@\finalhyphendemerits\z@
                \strut ##1\strut}%
        }%
        % If the linebreak is at a space, the latter will be displayed as visible
        % space at end of first line, and a continuation symbol starts next line.
        % Stretch/shrink are however usually zero for typewriter font.
        \def\FV@Space {%
            \nobreak\hskip\z@ plus\fontdimen3\font minus\fontdimen4\font
            \discretionary{\copy\Wrappedvisiblespacebox}{\Wrappedafterbreak}
            {\kern\fontdimen2\font}%
        }%

        % Allow breaks at special characters using \PYG... macros.
        \Wrappedbreaksatspecials
        % Breaks at punctuation characters . , ; ? ! and / need catcode=\active
        \OriginalVerbatim[#1,codes*=\Wrappedbreaksatpunct]%
    }
    \makeatother

    % Exact colors from NB
    \definecolor{incolor}{HTML}{303F9F}
    \definecolor{outcolor}{HTML}{D84315}
    \definecolor{cellborder}{HTML}{CFCFCF}
    \definecolor{cellbackground}{HTML}{F7F7F7}

    % prompt
    \makeatletter
    \newcommand{\boxspacing}{\kern\kvtcb@left@rule\kern\kvtcb@boxsep}
    \makeatother
    \newcommand{\prompt}[4]{
        {\ttfamily\llap{{\color{#2}[#3]:\hspace{3pt}#4}}\vspace{-\baselineskip}}
    }
    

    
    % Prevent overflowing lines due to hard-to-break entities
    \sloppy
    % Setup hyperref package
    \hypersetup{
      breaklinks=true,  % so long urls are correctly broken across lines
      colorlinks=true,
      urlcolor=urlcolor,
      linkcolor=linkcolor,
      citecolor=citecolor,
      }
    % Slightly bigger margins than the latex defaults
    
    \geometry{verbose,tmargin=1in,bmargin=1in,lmargin=1in,rmargin=1in}
    
    

\begin{document}
    
    \maketitle
    
    

    
    \section*{\texorpdfstring{\textbf{Anomaly vs.~Novelty Detection in
Medical
Data}}{Anomaly vs.~Novelty Detection in Medical Data}}\label{anomaly-vs.-novelty-detection-in-medical-data}

In modern healthcare, detecting unusual patterns in medical data is
crucial for early diagnosis and prevention of severe conditions. Machine
Learning (ML) offers powerful tools to identify abnormalities in patient
records, potentially improving disease detection and medical
decision-making.

In this project, we explore two key unsupervised learning techniques -
Anomaly Detection and Novelty Detection --- and apply them to a medical
dataset. While both techniques aim to identify unusual data points,
their applications differ:

\emph{Anomaly Detection} identifies outliers in a given dataset,
assuming that both normal and abnormal cases are present during
training. This method is useful for fraud detection, rare disease
identification, or malfunction detection in healthcare systems. \\
\emph{Novelty Detection} is trained only on normal cases and is designed
to detect previously unseen abnormalities. This is particularly useful
in early disease detection, where abnormal cases are rare or unavailable
during training.

For this study, we use a medical dataset (Breast Cancer) and implement
Isolation Forest and One-Class SVM models to compare their effectiveness
in distinguishing anomalies from normal patient records. The results
will provide insights into which technique is more effective for
detecting health-related abnormalities.

Through this analysis, we aim to: \\
1. Understand the differences between
Anomaly Detection and Novelty Detection.\\
2. Apply these techniques to
real-world medical data. \\
3. Compare their performance and practical
applications in healthcare.\\

This project contributes to the broader goal of leveraging machine
learning for improved medical diagnostics and anomaly detection in
patient records.

    \begin{tcolorbox}[breakable, size=fbox, boxrule=1pt, pad at break*=1mm,colback=cellbackground, colframe=cellborder]
\prompt{In}{incolor}{118}{\boxspacing}
\begin{Verbatim}[commandchars=\\\{\}]
\PY{c+c1}{\PYZsh{} Import necessary libraries}

\PY{k+kn}{import}\PY{+w}{ }\PY{n+nn}{numpy}\PY{+w}{ }\PY{k}{as}\PY{+w}{ }\PY{n+nn}{np}
\PY{k+kn}{import}\PY{+w}{ }\PY{n+nn}{pandas}\PY{+w}{ }\PY{k}{as}\PY{+w}{ }\PY{n+nn}{pd}
\PY{k+kn}{import}\PY{+w}{ }\PY{n+nn}{matplotlib}\PY{n+nn}{.}\PY{n+nn}{pyplot}\PY{+w}{ }\PY{k}{as}\PY{+w}{ }\PY{n+nn}{plt}
\PY{k+kn}{import}\PY{+w}{ }\PY{n+nn}{seaborn}\PY{+w}{ }\PY{k}{as}\PY{+w}{ }\PY{n+nn}{sns}

\PY{k+kn}{from}\PY{+w}{ }\PY{n+nn}{sklearn}\PY{n+nn}{.}\PY{n+nn}{model\PYZus{}selection}\PY{+w}{ }\PY{k+kn}{import} \PY{n}{train\PYZus{}test\PYZus{}split}
\PY{k+kn}{from}\PY{+w}{ }\PY{n+nn}{sklearn}\PY{n+nn}{.}\PY{n+nn}{preprocessing}\PY{+w}{ }\PY{k+kn}{import} \PY{n}{StandardScaler}
\PY{k+kn}{from}\PY{+w}{ }\PY{n+nn}{sklearn}\PY{n+nn}{.}\PY{n+nn}{ensemble}\PY{+w}{ }\PY{k+kn}{import} \PY{n}{IsolationForest}
\PY{k+kn}{from}\PY{+w}{ }\PY{n+nn}{sklearn}\PY{n+nn}{.}\PY{n+nn}{svm}\PY{+w}{ }\PY{k+kn}{import} \PY{n}{OneClassSVM}
\PY{k+kn}{from}\PY{+w}{ }\PY{n+nn}{sklearn}\PY{n+nn}{.}\PY{n+nn}{metrics}\PY{+w}{ }\PY{k+kn}{import} \PY{n}{classification\PYZus{}report}\PY{p}{,} \PY{n}{confusion\PYZus{}matrix}
\end{Verbatim}
\end{tcolorbox}

    \begin{tcolorbox}[breakable, size=fbox, boxrule=1pt, pad at break*=1mm,colback=cellbackground, colframe=cellborder]
\prompt{In}{incolor}{119}{\boxspacing}
\begin{Verbatim}[commandchars=\\\{\}]
\PY{k+kn}{from}\PY{+w}{ }\PY{n+nn}{ucimlrepo}\PY{+w}{ }\PY{k+kn}{import} \PY{n}{fetch\PYZus{}ucirepo} 
  
\PY{c+c1}{\PYZsh{} fetch dataset from online source}
\PY{n}{breast\PYZus{}cancer\PYZus{}dataset} \PY{o}{=} \PY{n}{fetch\PYZus{}ucirepo}\PY{p}{(}\PY{n+nb}{id}\PY{o}{=}\PY{l+m+mi}{17}\PY{p}{)} 

\PY{c+c1}{\PYZsh{} data (as pandas dataframes) }
\PY{n}{X} \PY{o}{=} \PY{n}{breast\PYZus{}cancer\PYZus{}dataset}\PY{o}{.}\PY{n}{data}\PY{o}{.}\PY{n}{features} 
\PY{n}{y} \PY{o}{=} \PY{n}{breast\PYZus{}cancer\PYZus{}dataset}\PY{o}{.}\PY{n}{data}\PY{o}{.}\PY{n}{targets} 

\PY{c+c1}{\PYZsh{} variable information }
\PY{n+nb}{print}\PY{p}{(}\PY{n}{breast\PYZus{}cancer\PYZus{}dataset}\PY{o}{.}\PY{n}{variables}\PY{p}{)} 
\end{Verbatim}
\end{tcolorbox}

    \begin{Verbatim}[commandchars=\\\{\}]
                  name     role         type demographic description units  \textbackslash{}
0                   ID       ID  Categorical        None        None  None
1            Diagnosis   Target  Categorical        None        None  None
2              radius1  Feature   Continuous        None        None  None
3             texture1  Feature   Continuous        None        None  None
4           perimeter1  Feature   Continuous        None        None  None
5                area1  Feature   Continuous        None        None  None
6          smoothness1  Feature   Continuous        None        None  None
7         compactness1  Feature   Continuous        None        None  None
8           concavity1  Feature   Continuous        None        None  None
9      concave\_points1  Feature   Continuous        None        None  None
10           symmetry1  Feature   Continuous        None        None  None
11  fractal\_dimension1  Feature   Continuous        None        None  None
12             radius2  Feature   Continuous        None        None  None
13            texture2  Feature   Continuous        None        None  None
14          perimeter2  Feature   Continuous        None        None  None
15               area2  Feature   Continuous        None        None  None
16         smoothness2  Feature   Continuous        None        None  None
17        compactness2  Feature   Continuous        None        None  None
18          concavity2  Feature   Continuous        None        None  None
19     concave\_points2  Feature   Continuous        None        None  None
20           symmetry2  Feature   Continuous        None        None  None
21  fractal\_dimension2  Feature   Continuous        None        None  None
22             radius3  Feature   Continuous        None        None  None
23            texture3  Feature   Continuous        None        None  None
24          perimeter3  Feature   Continuous        None        None  None
25               area3  Feature   Continuous        None        None  None
26         smoothness3  Feature   Continuous        None        None  None
27        compactness3  Feature   Continuous        None        None  None
28          concavity3  Feature   Continuous        None        None  None
29     concave\_points3  Feature   Continuous        None        None  None
30           symmetry3  Feature   Continuous        None        None  None
31  fractal\_dimension3  Feature   Continuous        None        None  None

   missing\_values
0              no
1              no
2              no
3              no
4              no
5              no
6              no
7              no
8              no
9              no
10             no
11             no
12             no
13             no
14             no
15             no
16             no
17             no
18             no
19             no
20             no
21             no
22             no
23             no
24             no
25             no
26             no
27             no
28             no
29             no
30             no
31             no
    \end{Verbatim}

    \begin{tcolorbox}[breakable, size=fbox, boxrule=1pt, pad at break*=1mm,colback=cellbackground, colframe=cellborder]
\prompt{In}{incolor}{120}{\boxspacing}
\begin{Verbatim}[commandchars=\\\{\}]
\PY{c+c1}{\PYZsh{}\PYZsh{} Viewing the first few rows}
\PY{n+nb}{print}\PY{p}{(}\PY{n}{X}\PY{o}{.}\PY{n}{head}\PY{p}{(}\PY{p}{)}\PY{p}{)}
\PY{c+c1}{\PYZsh{}\PYZsh{} Measures of central tendency}
\PY{n+nb}{print}\PY{p}{(}\PY{n}{X}\PY{o}{.}\PY{n}{describe}\PY{p}{(}\PY{p}{)}\PY{p}{)}
\end{Verbatim}
\end{tcolorbox}

    \begin{Verbatim}[commandchars=\\\{\}]
   radius1  texture1  perimeter1   area1  smoothness1  compactness1  \textbackslash{}
0    17.99     10.38      122.80  1001.0      0.11840       0.27760
1    20.57     17.77      132.90  1326.0      0.08474       0.07864
2    19.69     21.25      130.00  1203.0      0.10960       0.15990
3    11.42     20.38       77.58   386.1      0.14250       0.28390
4    20.29     14.34      135.10  1297.0      0.10030       0.13280

   concavity1  concave\_points1  symmetry1  fractal\_dimension1  {\ldots}  radius3  \textbackslash{}
0      0.3001          0.14710     0.2419             0.07871  {\ldots}    25.38
1      0.0869          0.07017     0.1812             0.05667  {\ldots}    24.99
2      0.1974          0.12790     0.2069             0.05999  {\ldots}    23.57
3      0.2414          0.10520     0.2597             0.09744  {\ldots}    14.91
4      0.1980          0.10430     0.1809             0.05883  {\ldots}    22.54

   texture3  perimeter3   area3  smoothness3  compactness3  concavity3  \textbackslash{}
0     17.33      184.60  2019.0       0.1622        0.6656      0.7119
1     23.41      158.80  1956.0       0.1238        0.1866      0.2416
2     25.53      152.50  1709.0       0.1444        0.4245      0.4504
3     26.50       98.87   567.7       0.2098        0.8663      0.6869
4     16.67      152.20  1575.0       0.1374        0.2050      0.4000

   concave\_points3  symmetry3  fractal\_dimension3
0           0.2654     0.4601             0.11890
1           0.1860     0.2750             0.08902
2           0.2430     0.3613             0.08758
3           0.2575     0.6638             0.17300
4           0.1625     0.2364             0.07678

[5 rows x 30 columns]
          radius1    texture1  perimeter1        area1  smoothness1  \textbackslash{}
count  569.000000  569.000000  569.000000   569.000000   569.000000
mean    14.127292   19.289649   91.969033   654.889104     0.096360
std      3.524049    4.301036   24.298981   351.914129     0.014064
min      6.981000    9.710000   43.790000   143.500000     0.052630
25\%     11.700000   16.170000   75.170000   420.300000     0.086370
50\%     13.370000   18.840000   86.240000   551.100000     0.095870
75\%     15.780000   21.800000  104.100000   782.700000     0.105300
max     28.110000   39.280000  188.500000  2501.000000     0.163400

       compactness1  concavity1  concave\_points1   symmetry1  \textbackslash{}
count    569.000000  569.000000       569.000000  569.000000
mean       0.104341    0.088799         0.048919    0.181162
std        0.052813    0.079720         0.038803    0.027414
min        0.019380    0.000000         0.000000    0.106000
25\%        0.064920    0.029560         0.020310    0.161900
50\%        0.092630    0.061540         0.033500    0.179200
75\%        0.130400    0.130700         0.074000    0.195700
max        0.345400    0.426800         0.201200    0.304000

       fractal\_dimension1  {\ldots}     radius3    texture3  perimeter3  \textbackslash{}
count          569.000000  {\ldots}  569.000000  569.000000  569.000000
mean             0.062798  {\ldots}   16.269190   25.677223  107.261213
std              0.007060  {\ldots}    4.833242    6.146258   33.602542
min              0.049960  {\ldots}    7.930000   12.020000   50.410000
25\%              0.057700  {\ldots}   13.010000   21.080000   84.110000
50\%              0.061540  {\ldots}   14.970000   25.410000   97.660000
75\%              0.066120  {\ldots}   18.790000   29.720000  125.400000
max              0.097440  {\ldots}   36.040000   49.540000  251.200000

             area3  smoothness3  compactness3  concavity3  concave\_points3  \textbackslash{}
count   569.000000   569.000000    569.000000  569.000000       569.000000
mean    880.583128     0.132369      0.254265    0.272188         0.114606
std     569.356993     0.022832      0.157336    0.208624         0.065732
min     185.200000     0.071170      0.027290    0.000000         0.000000
25\%     515.300000     0.116600      0.147200    0.114500         0.064930
50\%     686.500000     0.131300      0.211900    0.226700         0.099930
75\%    1084.000000     0.146000      0.339100    0.382900         0.161400
max    4254.000000     0.222600      1.058000    1.252000         0.291000

        symmetry3  fractal\_dimension3
count  569.000000          569.000000
mean     0.290076            0.083946
std      0.061867            0.018061
min      0.156500            0.055040
25\%      0.250400            0.071460
50\%      0.282200            0.080040
75\%      0.317900            0.092080
max      0.663800            0.207500

[8 rows x 30 columns]
    \end{Verbatim}

    \begin{tcolorbox}[breakable, size=fbox, boxrule=1pt, pad at break*=1mm,colback=cellbackground, colframe=cellborder]
\prompt{In}{incolor}{121}{\boxspacing}
\begin{Verbatim}[commandchars=\\\{\}]
\PY{c+c1}{\PYZsh{} Check for missing values}
\PY{n+nb}{print}\PY{p}{(}\PY{n}{X}\PY{o}{.}\PY{n}{isnull}\PY{p}{(}\PY{p}{)}\PY{o}{.}\PY{n}{sum}\PY{p}{(}\PY{p}{)}\PY{p}{)}
\end{Verbatim}
\end{tcolorbox}

    \begin{Verbatim}[commandchars=\\\{\}]
radius1               0
texture1              0
perimeter1            0
area1                 0
smoothness1           0
compactness1          0
concavity1            0
concave\_points1       0
symmetry1             0
fractal\_dimension1    0
radius2               0
texture2              0
perimeter2            0
area2                 0
smoothness2           0
compactness2          0
concavity2            0
concave\_points2       0
symmetry2             0
fractal\_dimension2    0
radius3               0
texture3              0
perimeter3            0
area3                 0
smoothness3           0
compactness3          0
concavity3            0
concave\_points3       0
symmetry3             0
fractal\_dimension3    0
dtype: int64
    \end{Verbatim}

    \section{Dataset}\label{dataset}

The breast cancer dataset has 569 rows of data with 32 attributes (ID,
Diagnosis and 30 real-valued input features).

The Diagnosis (target) is either M = malignant and B = benign. Ten
real-valued features are computed fro each cell nucleus;

\begin{verbatim}
a) radius (mean of distances from center to points on the perimeter)
b) texture (standard deviation of gray-scale values)
c) perimeter
d) area
e) smoothness (local variation in radius lengths)
f) compactness (perimeter^2 / area - 1.0)
g) concavity (severity of concave portions of the contour)
h) concave points (number of concave portions of the contour)
i) symmetry 
j) fractal dimension ("coastline approximation" - 1)
\end{verbatim}

The dataset has no missing values.

    \begin{tcolorbox}[breakable, size=fbox, boxrule=1pt, pad at break*=1mm,colback=cellbackground, colframe=cellborder]
\prompt{In}{incolor}{122}{\boxspacing}
\begin{Verbatim}[commandchars=\\\{\}]
\PY{c+c1}{\PYZsh{}Standadize the features }

\PY{n}{scaler} \PY{o}{=} \PY{n}{StandardScaler}\PY{p}{(}\PY{p}{)}
\PY{n}{X\PYZus{}scaled} \PY{o}{=} \PY{n}{scaler}\PY{o}{.}\PY{n}{fit\PYZus{}transform}\PY{p}{(}\PY{n}{X}\PY{p}{)}
\end{Verbatim}
\end{tcolorbox}

    \begin{tcolorbox}[breakable, size=fbox, boxrule=1pt, pad at break*=1mm,colback=cellbackground, colframe=cellborder]
\prompt{In}{incolor}{123}{\boxspacing}
\begin{Verbatim}[commandchars=\\\{\}]
\PY{c+c1}{\PYZsh{} Splitting the data}
\PY{n}{X\PYZus{}train}\PY{p}{,} \PY{n}{X\PYZus{}test}\PY{p}{,} \PY{n}{y\PYZus{}train}\PY{p}{,} \PY{n}{y\PYZus{}test} \PY{o}{=} \PY{n}{train\PYZus{}test\PYZus{}split}\PY{p}{(}\PY{n}{X\PYZus{}scaled}\PY{p}{,} \PY{n}{y}\PY{p}{,} \PY{n}{test\PYZus{}size}\PY{o}{=}\PY{l+m+mf}{0.2}\PY{p}{,} \PY{n}{random\PYZus{}state}\PY{o}{=}\PY{l+m+mi}{42}\PY{p}{)}
\end{Verbatim}
\end{tcolorbox}

    \section{Anomaly Detection with Isolation
forest}\label{anomaly-detection-with-isolation-forest}

\begin{itemize}
\tightlist
\item
  Isolation Forest is a tree-based algorithm that identifies anomalies
  by splitting data points recursively.
\item
  It isolates outliers faster than normal points because they require
  fewer splits.
\item
  It assigns a score to each point, and if a point is far from the main
  distribution, it is labeled as an anomaly.
\end{itemize}

\subsection{Isolation Forest Model}\label{isolation-forest-model}

To initialize the isolation forest model, we should define the number of
base estimators or trees in the ensemble (\emph{n\_estimators}), the
\emph{contamination} parameter, which assumes the \% of the data that
are anomalies and the \emph{random\_state} which ensures results are
reproducible.

The .predict function assigns labels -1 for Anomaly and 1 for Normal

    \begin{tcolorbox}[breakable, size=fbox, boxrule=1pt, pad at break*=1mm,colback=cellbackground, colframe=cellborder]
\prompt{In}{incolor}{124}{\boxspacing}
\begin{Verbatim}[commandchars=\\\{\}]
\PY{c+c1}{\PYZsh{} Train Isolation Forest model}
\PY{n}{iso\PYZus{}forest} \PY{o}{=} \PY{n}{IsolationForest}\PY{p}{(}\PY{n}{n\PYZus{}estimators}\PY{o}{=}\PY{l+m+mi}{100}\PY{p}{,} \PY{n}{contamination}\PY{o}{=}\PY{l+m+mf}{0.1}\PY{p}{,} \PY{n}{random\PYZus{}state}\PY{o}{=}\PY{l+m+mi}{42}\PY{p}{)}
\PY{n}{iso\PYZus{}forest}\PY{o}{.}\PY{n}{fit}\PY{p}{(}\PY{n}{X\PYZus{}train}\PY{p}{)}

\PY{c+c1}{\PYZsh{} Predict anomalies (\PYZhy{}1 = Anomaly, 1 = Normal)}
\PY{n}{y\PYZus{}pred\PYZus{}iso} \PY{o}{=} \PY{n}{iso\PYZus{}forest}\PY{o}{.}\PY{n}{predict}\PY{p}{(}\PY{n}{X\PYZus{}test}\PY{p}{)}

\PY{c+c1}{\PYZsh{} Convert predictions to 0 (normal) and 1 (anomaly)}
\PY{n}{y\PYZus{}test} \PY{o}{=} \PY{n}{np}\PY{o}{.}\PY{n}{where}\PY{p}{(} \PY{n}{y\PYZus{}test} \PY{o}{==} \PY{l+s+s1}{\PYZsq{}}\PY{l+s+s1}{M}\PY{l+s+s1}{\PYZsq{}}\PY{p}{,} \PY{l+m+mi}{1}\PY{p}{,} \PY{l+m+mi}{0}\PY{p}{)}
\PY{n}{y\PYZus{}pred\PYZus{}iso} \PY{o}{=} \PY{n}{np}\PY{o}{.}\PY{n}{where}\PY{p}{(}\PY{n}{y\PYZus{}pred\PYZus{}iso} \PY{o}{==} \PY{o}{\PYZhy{}}\PY{l+m+mi}{1}\PY{p}{,} \PY{l+m+mi}{1}\PY{p}{,} \PY{l+m+mi}{0}\PY{p}{)}
\end{Verbatim}
\end{tcolorbox}

    \begin{tcolorbox}[breakable, size=fbox, boxrule=1pt, pad at break*=1mm,colback=cellbackground, colframe=cellborder]
\prompt{In}{incolor}{125}{\boxspacing}
\begin{Verbatim}[commandchars=\\\{\}]
\PY{c+c1}{\PYZsh{} Evaluate the model}
\PY{n+nb}{print}\PY{p}{(}\PY{l+s+s2}{\PYZdq{}}\PY{l+s+s2}{Isolation Forest Classification Report:}\PY{l+s+s2}{\PYZdq{}}\PY{p}{)}
\PY{n+nb}{print}\PY{p}{(}\PY{n}{classification\PYZus{}report}\PY{p}{(}\PY{n}{y\PYZus{}test}\PY{p}{,} \PY{n}{y\PYZus{}pred\PYZus{}iso}\PY{p}{)}\PY{p}{)}

\PY{c+c1}{\PYZsh{} Confusion matrix}
\PY{n}{sns}\PY{o}{.}\PY{n}{heatmap}\PY{p}{(}\PY{n}{confusion\PYZus{}matrix}\PY{p}{(}\PY{n}{y\PYZus{}test}\PY{p}{,} \PY{n}{y\PYZus{}pred\PYZus{}iso}\PY{p}{)}\PY{p}{,} \PY{n}{annot}\PY{o}{=}\PY{k+kc}{True}\PY{p}{,} \PY{n}{fmt}\PY{o}{=}\PY{l+s+s1}{\PYZsq{}}\PY{l+s+s1}{d}\PY{l+s+s1}{\PYZsq{}}\PY{p}{,} \PY{n}{cmap}\PY{o}{=}\PY{l+s+s2}{\PYZdq{}}\PY{l+s+s2}{Reds}\PY{l+s+s2}{\PYZdq{}}\PY{p}{)}
\PY{n}{plt}\PY{o}{.}\PY{n}{xlabel}\PY{p}{(}\PY{l+s+s2}{\PYZdq{}}\PY{l+s+s2}{Predicted}\PY{l+s+s2}{\PYZdq{}}\PY{p}{)}
\PY{n}{plt}\PY{o}{.}\PY{n}{ylabel}\PY{p}{(}\PY{l+s+s2}{\PYZdq{}}\PY{l+s+s2}{Actual}\PY{l+s+s2}{\PYZdq{}}\PY{p}{)}
\PY{n}{plt}\PY{o}{.}\PY{n}{title}\PY{p}{(}\PY{l+s+s2}{\PYZdq{}}\PY{l+s+s2}{Confusion Matrix \PYZhy{} Isolation Forest}\PY{l+s+s2}{\PYZdq{}}\PY{p}{)}
\PY{n}{plt}\PY{o}{.}\PY{n}{show}\PY{p}{(}\PY{p}{)}
\end{Verbatim}
\end{tcolorbox}

    \begin{Verbatim}[commandchars=\\\{\}]
Isolation Forest Classification Report:
              precision    recall  f1-score   support

           0       0.68      0.94      0.79        71
           1       0.73      0.26      0.38        43

    accuracy                           0.68       114
   macro avg       0.71      0.60      0.58       114
weighted avg       0.70      0.68      0.63       114

    \end{Verbatim}

    \begin{center}
    \adjustimage{max size={0.9\linewidth}{0.9\paperheight}}{output_10_1.png}
    \end{center}
    { \hspace*{\fill} \\}
    
    \subsection{Interpretation of Results}\label{interpretation-of-results}

\begin{itemize}
\tightlist
\item
  \textbf{Precision = 0.73} → When the model predicts an anomaly, it is
  correct 73\% of the time.
\item
  \textbf{Recall = 0.26} → The model is only detecting 26\% of actual
  anomalies, meaning it's missing 74\% of anomalies (high false
  negatives).
\item
  \textbf{F1-score = 0.38} → The balance between precision and recall is
  low, suggesting a weak overall model.
\item
  \textbf{Accuracy = 0.68} → The model correctly classifies 68\% of all
  cases, but since anomalies are a minority, accuracy can be misleading.
\end{itemize}

Some potential issues would be; 1. The model is too conservative meaning
it only marks cases as anomalies when it's very confident, leading to
many missed anomalies (low recall). 2. Anomalies are underrepresented in
training making the model struggle to learn them.

The model can be improved by;

\begin{enumerate}
\def\labelenumi{\arabic{enumi}.}
\tightlist
\item
  Increase Recall (detect more anomalies):

  \begin{itemize}
  \tightlist
  \item
    Increase \emph{contamination} to allow more anomalies.
  \item
    Try a different anomaly detection method, such as LOF (Local Outlier
    Factor) or Autoencoders.
  \end{itemize}
\item
  Balance Precision \& Recall (improve F1-score):

  \begin{itemize}
  \tightlist
  \item
    Try an ensemble approach, combining multiple anomaly detection
    models.
  \item
    Use feature engineering to improve the separability of normal
    vs.~anomalous cases.
  \end{itemize}
\end{enumerate}

    \section{Novelty Detection with One-Class
SVM}\label{novelty-detection-with-one-class-svm}

\begin{itemize}
\tightlist
\item
  One-Class SVM (Support Vector Machine) is a boundary-based model. It
  learns the shape of normal data and flags anything outside this
  boundary as a novel case.
\item
  Difference from Isolation Forest: Here, we assume training data
  contains only normal cases. It does not learn from existing anomalies,
  so it's useful for detecting previously unseen abnormalities.
\end{itemize}

\subsection{One-Class SVM}\label{one-class-svm}

Since novelty detection assumes we have only normal training data, we
first detect and filter out anomalies using isolation forest. Train
One-Class SVM on only normal cases. \emph{nu}=0.05 means the model
allows 5\% deviation from normal whereas \emph{kernel=``rbf''} ensures
smooth decision boundaries.

Like Isolation Forest: -1 represents Novel (previously unseen) case and
1 represents Normal case If the model detects cases that weren't seen
during training, they will be flagged.

    \begin{tcolorbox}[breakable, size=fbox, boxrule=1pt, pad at break*=1mm,colback=cellbackground, colframe=cellborder]
\prompt{In}{incolor}{126}{\boxspacing}
\begin{Verbatim}[commandchars=\\\{\}]
\PY{c+c1}{\PYZsh{} Step 1: Use Isolation Forest to detect anomalies in the training set}
\PY{n}{iso\PYZus{}forest} \PY{o}{=} \PY{n}{IsolationForest}\PY{p}{(}\PY{n}{contamination}\PY{o}{=}\PY{l+m+mf}{0.05}\PY{p}{,} \PY{n}{random\PYZus{}state}\PY{o}{=}\PY{l+m+mi}{42}\PY{p}{)}  \PY{c+c1}{\PYZsh{} Adjust contamination as needed}
\PY{n}{outlier\PYZus{}labels} \PY{o}{=} \PY{n}{iso\PYZus{}forest}\PY{o}{.}\PY{n}{fit\PYZus{}predict}\PY{p}{(}\PY{n}{X\PYZus{}scaled}\PY{p}{)}

\PY{c+c1}{\PYZsh{} Step 3: Filter out anomalies (keep only normal data for training)}
\PY{n}{X\PYZus{}train\PYZus{}cleaned} \PY{o}{=} \PY{n}{X\PYZus{}scaled}\PY{p}{[}\PY{n}{outlier\PYZus{}labels} \PY{o}{==} \PY{l+m+mi}{1}\PY{p}{]}  \PY{c+c1}{\PYZsh{} 1 = normal, \PYZhy{}1 = anomaly}

\PY{c+c1}{\PYZsh{} Step 4: Train One\PYZhy{}Class SVM on the cleaned training data}
\PY{n}{oc\PYZus{}svm} \PY{o}{=} \PY{n}{OneClassSVM}\PY{p}{(}\PY{n}{kernel}\PY{o}{=}\PY{l+s+s2}{\PYZdq{}}\PY{l+s+s2}{rbf}\PY{l+s+s2}{\PYZdq{}}\PY{p}{,} \PY{n}{nu}\PY{o}{=}\PY{l+m+mf}{0.05}\PY{p}{,} \PY{n}{gamma}\PY{o}{=}\PY{l+s+s2}{\PYZdq{}}\PY{l+s+s2}{scale}\PY{l+s+s2}{\PYZdq{}}\PY{p}{)}
\PY{n}{oc\PYZus{}svm}\PY{o}{.}\PY{n}{fit}\PY{p}{(}\PY{n}{X\PYZus{}train\PYZus{}cleaned}\PY{p}{)}

\PY{c+c1}{\PYZsh{} Step 5: Predict anomalies on the test set (using the full dataset)}
\PY{n}{predictions} \PY{o}{=} \PY{n}{oc\PYZus{}svm}\PY{o}{.}\PY{n}{predict}\PY{p}{(}\PY{n}{X\PYZus{}scaled}\PY{p}{)}

\PY{c+c1}{\PYZsh{} Convert predictions: 1 \PYZhy{}\PYZgt{} Normal, \PYZhy{}1 \PYZhy{}\PYZgt{} Anomaly}
\PY{n}{predictions} \PY{o}{=} \PY{n}{np}\PY{o}{.}\PY{n}{where}\PY{p}{(}\PY{n}{predictions} \PY{o}{==} \PY{l+m+mi}{1}\PY{p}{,} \PY{l+m+mi}{1}\PY{p}{,} \PY{l+m+mi}{0}\PY{p}{)}
\PY{n}{y} \PY{o}{=} \PY{n}{np}\PY{o}{.}\PY{n}{where}\PY{p}{(} \PY{n}{y} \PY{o}{==} \PY{l+s+s1}{\PYZsq{}}\PY{l+s+s1}{M}\PY{l+s+s1}{\PYZsq{}}\PY{p}{,} \PY{l+m+mi}{1}\PY{p}{,} \PY{l+m+mi}{0}\PY{p}{)}

\PY{c+c1}{\PYZsh{} Evaluate the model}
\PY{n+nb}{print}\PY{p}{(}\PY{l+s+s2}{\PYZdq{}}\PY{l+s+s2}{One\PYZhy{}Class SVM Classification Report:}\PY{l+s+s2}{\PYZdq{}}\PY{p}{)}
\PY{n+nb}{print}\PY{p}{(}\PY{n}{classification\PYZus{}report}\PY{p}{(}\PY{n}{y}\PY{p}{,} \PY{n}{predictions}\PY{p}{)}\PY{p}{)}

\PY{c+c1}{\PYZsh{} Confusion matrix}
\PY{n}{sns}\PY{o}{.}\PY{n}{heatmap}\PY{p}{(}\PY{n}{confusion\PYZus{}matrix}\PY{p}{(}\PY{n}{y}\PY{p}{,} \PY{n}{predictions}\PY{p}{)}\PY{p}{,} \PY{n}{annot}\PY{o}{=}\PY{k+kc}{True}\PY{p}{,} \PY{n}{fmt}\PY{o}{=}\PY{l+s+s1}{\PYZsq{}}\PY{l+s+s1}{d}\PY{l+s+s1}{\PYZsq{}}\PY{p}{,} \PY{n}{cmap}\PY{o}{=}\PY{l+s+s2}{\PYZdq{}}\PY{l+s+s2}{Greens}\PY{l+s+s2}{\PYZdq{}}\PY{p}{)}
\PY{n}{plt}\PY{o}{.}\PY{n}{xlabel}\PY{p}{(}\PY{l+s+s2}{\PYZdq{}}\PY{l+s+s2}{Predicted}\PY{l+s+s2}{\PYZdq{}}\PY{p}{)}
\PY{n}{plt}\PY{o}{.}\PY{n}{ylabel}\PY{p}{(}\PY{l+s+s2}{\PYZdq{}}\PY{l+s+s2}{Actual}\PY{l+s+s2}{\PYZdq{}}\PY{p}{)}
\PY{n}{plt}\PY{o}{.}\PY{n}{title}\PY{p}{(}\PY{l+s+s2}{\PYZdq{}}\PY{l+s+s2}{Confusion Matrix \PYZhy{} One\PYZhy{}Class SVM}\PY{l+s+s2}{\PYZdq{}}\PY{p}{)}
\PY{n}{plt}\PY{o}{.}\PY{n}{show}\PY{p}{(}\PY{p}{)}
\end{Verbatim}
\end{tcolorbox}

    \begin{Verbatim}[commandchars=\\\{\}]
One-Class SVM Classification Report:
              precision    recall  f1-score   support

           0       0.45      0.08      0.14       357
           1       0.35      0.83      0.49       212

    accuracy                           0.36       569
   macro avg       0.40      0.46      0.31       569
weighted avg       0.41      0.36      0.27       569

    \end{Verbatim}

    \begin{center}
    \adjustimage{max size={0.9\linewidth}{0.9\paperheight}}{output_13_1.png}
    \end{center}
    { \hspace*{\fill} \\}
    
    \subsection{Interpretation of Results}\label{interpretation-of-results}

\begin{itemize}
\tightlist
\item
  \textbf{Precision = 0.35} → This suggests that only 35\% of the points
  the model flagged as anomalies were truly anomalies. This is
  relatively low meaning the model is incorrectly labelling many normal
  cases as anomalies (high false positives)
\item
  \textbf{Recall = 0.83} → This means 83\% of the actual anomalies were
  correctly detected.
\item
  \textbf{F1-score = 0.49} → An F1-score of 0.49 suggests that the model
  is somewhat unbalanced, with a high recall but low precision.
\item
  \textbf{Accuracy = 0.36} → An accuracy of 0.36 is quite low,
  suggesting that the model might be over-predicting anomalies or
  failing to properly classify normal instances, especially if the
  dataset is imbalanced.
\end{itemize}

Recommendations to improve the model could be;

\begin{enumerate}
\def\labelenumi{\arabic{enumi}.}
\tightlist
\item
  Adjust the \emph{nu} parameter in One-Class SVM:

  \begin{itemize}
  \tightlist
  \item
    The \emph{nu} parameter controls the upper bound on the fraction of
    margin errors (false positives). A high value might be leading to
    more false positives, so try lowering it to see if precision
    improves.
  \end{itemize}
\item
  Consider tuning the gamma parameter:

  \begin{itemize}
  \tightlist
  \item
    This controls the influence of individual data points. A large gamma
    can lead to overfitting, while a small gamma may underfit. Adjusting
    it can help improve model performance.
  \end{itemize}
\item
  Try other anomaly detection methods:

  \begin{itemize}
  \tightlist
  \item
    If One-Class SVM isn't giving satisfactory results, consider trying
    methods like Isolation Forest or Autoencoders for anomaly detection.
    They might provide a better fit depending on your data.Use feature
    engineering to improve the separability of normal vs.~anomalous
    cases.
  \end{itemize}
\end{enumerate}

    \section{Conclusion}\label{conclusion}

\subsection{Objective Recap:}\label{objective-recap}

The objective of this project was to explore and compare anomaly
detection and novelty detection techniques using the Breast Cancer
Dataset. Specifically, we aimed to identify and classify anomalous
(malignant) tumors from a larger set of normal (benign) tumor data using
One-Class SVM for novelty detection and Isolation Forest for anomaly
detection.

\subsection{Key Insights and Results:}\label{key-insights-and-results}

\subsubsection{Anomaly Detection with Isolation
Forest:}\label{anomaly-detection-with-isolation-forest}

Isolation Forest proved effective in identifying anomalies (malignant
tumors), but the model showed some limitations in recall and precision.
The model was able to detect a good number of anomalies (with relatively
high recall) but exhibited a moderate precision due to misclassifying
some normal cases as anomalies. The accuracy was higher, indicating it
could generally distinguish between normal and anomalous data, but some
improvements could be made in capturing all anomalies.

\subsubsection{Novelty Detection with One-Class
SVM:}\label{novelty-detection-with-one-class-svm}

One-Class SVM achieved high recall (83\%), indicating its strength in
detecting true anomalies, but low precision (35\%), suggesting many
normal cases were wrongly identified as anomalies (false positives). The
F1-score and accuracy were lower, revealing that despite its ability to
detect most anomalies, the model suffers from a high rate of false
positives, impacting its reliability for practical applications.

\subsection{Comparison of Techniques:}\label{comparison-of-techniques}

One-Class SVM performed well in terms of recall, which is often more
critical in anomaly detection tasks where catching all outliers is a
priority, but its low precision makes it less reliable for applications
where false positives should be minimized. Isolation Forest offered a
balanced approach with higher accuracy, though further adjustments in
hyperparameters could improve both precision and recall.

\subsection{Challenges and Areas for
Improvement:}\label{challenges-and-areas-for-improvement}

\begin{itemize}
\tightlist
\item
  \emph{Model Tuning}: Both models could benefit from more
  hyperparameter tuning (e.g., adjusting nu in One-Class SVM and
  contamination in Isolation Forest) to improve precision and reduce
  false positives.
\item
  \emph{Data Imbalance}: The dataset's potential imbalance between
  normal and anomalous cases might have skewed the results, particularly
  affecting the precision and accuracy. Techniques like SMOTE or class
  weighting could improve model performance.
\item
  \emph{Feature Engineering}: Further feature engineering or exploration
  of other input features could help enhance model performance,
  especially in capturing subtle patterns in the data.
\end{itemize}

\subsection{Key Takeaways and Next
Steps:}\label{key-takeaways-and-next-steps}

Anomaly and novelty detection methods are powerful tools for identifying
outliers in data, particularly when the data is imbalanced. However,
fine-tuning model parameters is crucial for improving precision and
recall. For real-world applications in health data or fraud detection,
balancing precision and recall is essential to ensure the model is both
accurate and reliable. Moving forward, we could explore more
sophisticated models like Autoencoders for anomaly detection or even
ensemble methods for improved robustness.

\subsection{Final Thoughts:}\label{final-thoughts}

This project has provided valuable insights into the application of
anomaly and novelty detection techniques. While challenges remain in
optimizing models for precision and recall, the experience gained here
offers a solid foundation for tackling real-world anomaly detection
problems, especially in fields like medical diagnostics and fraud
detection.


    % Add a bibliography block to the postdoc
    
    
    
\end{document}
